% Revised Introduction - Publication Quality
% Tells a clear story: Problem → Gap → Approach → Contributions

\section{Introduction}

\subsection{Motivation: The Observation Process Problem}

Industrial prognostic systems operate in environments where sensor configurations rarely remain static. Communication dropouts create intermittent data gaps. Individual sensors fail and await replacement. Maintenance interventions temporarily disable portions of the sensing infrastructure. During these events, predictive maintenance models must continue to provide reliable remaining useful life (RUL) estimates, yet most deep learning architectures assume the observation process---the subset and quality of available sensors---remains fixed throughout deployment.

This mismatch creates a subtle but consequential failure mode. A model may achieve acceptable pointwise prediction accuracy while producing \emph{incompatible posterior distributions} when the same asset is viewed through different observation configurations. For example, predictions made with a full sensor complement may yield narrow, confident distributions, while predictions from a reduced sensor set produce wider distributions that do not properly overlap the full-sensor posterior---even though both views ostensibly describe the same underlying degradation state. Such geometric inconsistencies undermine uncertainty-aware decision making, where maintenance scheduling depends not just on the posterior mean but on tail probabilities, prediction intervals, and risk assessments.

\subsection{Prior Work and the Unresolved Question}

The methodological foundation separates the observation mechanism from the latent
quantity being predicted~\citep{rubin1976inference,little2019statistical}. Neural
imputation methods such as BRITS and CSDI~\citep{cao2018brits,tashiro2021csdi},
permutation-invariant set encoders and attention mechanisms~\citep{lee2019settransformer,vaswani2017attention,zaheer2017deep},
and conditional neural processes~\citep{garnelo2018cnp} each provide a way to
represent incomplete or changing inputs. For probabilistic prognostics, proper
scoring rules and calibration diagnostics provide complementary checks on forecast
quality~\citep{gneiting2007scoring,kuleshov2018calibrated,angelopoulos202300000101}.
These foundations motivate the observation-process perspective, while the RUL
literature below supplies the closest application-specific comparisons.

Recent prognostics research has begun to address variable and degraded sensing, but
its primary objectives differ across several neighboring lines of work. Methods for
missing or faulty sensors typically aim to preserve point-prediction accuracy by
recovering the input or learning a representation that is robust to incomplete
channels. Niu et al.~\citep{niu2026dynamic} use real-time data imputation within a
dynamic adaptive prognostics pipeline, while Liu et al.~\citep{liu2026strsmtnet}
jointly optimize sensor reconstruction and RUL prediction with a spatio-temporal
multi-task network. Zhang et al.~\citep{zhang2025msccan} introduce a variable-sensor
scenario and use cross-channel attention and data augmentation to stabilize RUL
prediction. Wang et al.~\citep{wang2025gfggat} model partially observed sensor
streams with a graph feature-gated graph attention network,
thereby making relations among sensor channels explicit. These studies establish
strong precedents for learning under changing sensor availability, yet their main
comparisons remain single-view prediction errors against the observed RUL target.

A second line of work focuses on uncertainty quantification and reliability. Sankararaman~\citep{sankararaman201501405029} frames prognostic uncertainty as an uncertainty
propagation problem, and Li et al.~\citep{li202103009593} distinguish epistemic and
aleatoric contributions in a Bayesian deep-learning RUL framework. More recent work
combines missing-value modeling with distribution-free interval construction:
Wang et al.~\citep{wang2025robustuq} use diffusion-based imputation, metric learning,
and conformal prediction to obtain statistically valid intervals under randomly
missing or partially faulty sensors. Xu et al.~\citep{xu2026calibratedrul} further
combine Bayesian predictive distributions with isotonic uncertainty calibration and
split conformal intervals for RUL prediction. The recent review by Lin et al.~\citep{lin2026uncertaintyreview}
places these approaches within a broader taxonomy of statistical, learning-based,
and hybrid uncertainty analysis. Related work on conformal RUL intervals and domain
generalization~\citep{javanmardi202314i23417,xu202310381131,ding202322110199,tong202443366689}
likewise emphasizes coverage, robustness, or transfer across operating conditions.

Sensor degradation extends the same concern beyond binary availability. Earlier
reliability studies modelled joint system--sensor degradation in RUL inference using
state-space and filtering methods~\citep{mukhopadhyay2023degradingsensors,hachem2024sensordegradation}.
Zhao et al.~\citep{zhao2026sensordegradation} extend this line to rolling bearings by
jointly modelling bearing and sensor degradation so that measurement quality can be
updated during RUL inference. This perspective is useful for defining the observation
process more broadly: availability, measurement quality, and functional sensor
importance are related but distinct aspects of what a predictor observes.

These neighboring methods answer important questions about recovering missing
measurements, maintaining point accuracy, constructing valid intervals, or exploiting
sensor relations. However, they do not make the paired geometric relationship between
predictive distributions from different sensor views of the same asset their primary
estimand. The question studied here is different: for the same asset and query time,
when two sensor sets remain task-sufficient, do they induce posterior distributions
with a coherent geometric relationship? We therefore treat the observation process
as an explicit conditioning variable and make cross-view posterior inconsistency the
primary estimand. This distinction also explains the role of the boundary analysis:
a model may recognize that a sensor set has changed while still failing to represent
the functional importance of the missing channels.

Time-series foundation models have broadened the forecasting model family and
motivated more realistic transfer analyses~\citep{rasul202331008278,ansari202440307815,qiu202443665863}.
They provide useful modeling backbones, but do not by themselves define a paired
probabilistic estimand for changing sensor sets. Our study isolates that
observation-process question within a controlled prognostic setting.

\subsection{Our Approach: Explicit Observation-Process Conditioning}

We propose an architecture that treats the observation process as a \emph{first-class conditioning variable}, analogous to how conditional generative models explicitly condition on class labels or attributes. The key insight is that posterior geometry should \emph{explicitly} depend on which sensors are available, not merely adapt implicitly through input masking.

Concretely, our model---termed Prognostic Observation-Process-Directed (PROP-DIRECT)---embeds the sensor set configuration into a latent representation, then conditions all posterior computations on this representation alongside the observed sensor values. This design provides three potential benefits:

\begin{enumerate}
    \item \textbf{Explicit uncertainty propagation:} The model learns to widen posteriors appropriately when critical sensors are missing, rather than maintaining false confidence.
    \item \textbf{Cross-view consistency:} Predictions from different observation views of the same asset should remain geometrically compatible---overlapping posteriors for overlapping information.
    \item \textbf{Graceful degradation:} As sensor availability decreases, posterior uncertainty should increase smoothly rather than exhibiting discontinuous failures.
\end{enumerate}

\subsection{Scope and Contributions}

We test this question with a preregistered protocol that separates model development from independent test and lockbox evaluations. The central hypothesis is conditional: conditioning should help when alternative sensor views remain \emph{task-sufficient}, meaning that the available measurements still support reliable RUL prediction. The protocol defines task sufficiency through oracle-posterior quality thresholds (Section~\ref{sec:methods:protocol}) and then examines what happens when functionally critical sensors are removed.

Our contributions are:

\begin{enumerate}
    \item \textbf{Empirical result under task sufficiency:} Across 96 synthetic assets with controlled observation processes, sensor-set conditioning reduces cross-view posterior inconsistency by 38.6\% (90\% CI: 34.8--42.4\%) relative to observation-agnostic baselines. A complementary shared-model analysis on the public C-MAPSS simulated benchmark reveals that the geometry gain depends on which view pair is compared (Section~\ref{sec:results_cmapss}).

    \item \textbf{Boundary characterization under information loss:} When critical sensors are removed, the architecture does not maintain its advantage across three pre-registered superiority criteria, locating the boundary of the conditioning mechanism.

    \item \textbf{Mechanistic insight via exploratory analysis:} Post-hoc correlation analysis reveals \emph{why} the boundary exists: the model successfully detects magnitude of information loss (oracle-model correlation $r>0.53$), but shows a 9.9\% weaker correlation when importance matters. Because prediction cases share assets and query structure, this contrast is treated as exploratory rather than as a confirmatory independent-sample test. Crucially, the pattern is architecture-specific---a standard baseline shows the opposite direction.

    \item \textbf{Quality-geometry separation:} In the synthetic information-loss evaluation, predictive quality remains competitive even when geometric consistency degrades, separating posterior shape from point prediction accuracy.
\end{enumerate}

\subsection{Study Design and Scope}

We use a synthetic generator because exact oracle posteriors and controlled observation processes make the geometric comparison identifiable. A complementary analysis on the public C-MAPSS simulated turbofan benchmark asks whether the scoring pattern extends beyond this generator; the Discussion places that result alongside the still-open question of naturally occurring sensor loss.

\subsection{Paper Organization}

Section~\ref{sec:methods} describes the PROP-DIRECT architecture, study protocol, metrics, and benchmark design. Section~\ref{sec:results} reports the primary analysis under sufficient observation, the information-loss boundary, the exploratory mechanism analysis, predictive quality and calibration, and the complementary C-MAPSS evaluation. Section~\ref{sec:discussion} interprets these findings, positions them relative to existing approaches, and provides deployment guidance. Section~\ref{sec:conclusion} summarizes contributions and outlines future directions.
